\chapter{Conclusiones}

Las conclusiones definitivas se redactarán una vez completada la evaluación experimental. De forma provisional, el análisis desarrollado hasta el momento en el TFE indica que la búsqueda cuántica de máximo puede reducir el número de consultas sobre el espacio de subclaves, mientras que la puntuación continúa dependiendo del volumen de datos y de un oráculo reversible cuyo coste debe contabilizarse. La magnitud de la ventaja total no puede establecerse hasta fijar y reproducir la instancia concreta.

El trabajo futuro inmediato consiste en inspeccionar y documentar el código fuente, verificar la función de conteo frente a casos pequeños, ejecutar experimentos con parámetros registrados y construir el circuito reversible correspondiente. La evaluación final deberá considerar también la carga coherente de datos, profundidad, número de qubits, precisión de la búsqueda y supuestos de acceso al cifrador. Estas tareas permitirán distinguir una mejora de complejidad en consultas de una aceleración práctica.
